\section{Introduction}
\label{sec:intro}

Deployed language-model agents are specialists. A coding agent depends
on code synthesis and tool use, a retrieval assistant on multi-hop
question answering, and neither needs frontier-scale generality while
both pay for every token of inference. The practical question is which
model covers a workload's capability profile at the lowest cost, and
the usual answer starts from a strong dense source and scales it down
by pruning, quantization, or distillation, often with recovery
training. Compression does not shrink capabilities uniformly. At the
same size or bit-width, different methods and different source
families preserve different capabilities, improve some while collapsing
others, and hit cliffs at different strengths (\S\ref{sec:phenomena}).
Aggregate perplexity and average benchmark scores hide this, so
choosing a route today means running each candidate configuration
through capability benchmarks.

This paper asks whether that search can be predicted from what is
observable before compression. We fix a per-capability conditional
loss as the endpoint and write every prediction relative to a reference
measured beforehand: the dense source for pruning and quantization, the
initial student for distillation (\S\ref{sec:prelim}). The admitted
inputs are the source size, its pretraining token count where
disclosed, its dense capability losses, and the intervention setting.
The question is then concrete. Can pre-compression model-state
information and intervention settings predict capability-conditioned
loss, and over which source and intervention ranges do those
predictions transfer?

Answering it requires separating three things that a single accuracy
number would blur. A predictor can fail because the source information
is useless, because the functional form is wrong, or because the test
lies outside the range in which anything was fit. We therefore test
two axes separately, transfer across source-model states at a fixed
intervention and prediction at unseen intervention settings, and we
compare every candidate against baselines that use the same
information: a source-free strength curve or a per-configuration
constant asks whether the source state helps at all, a no-$D_0$ variant
asks what training tokens add, and same-input alternatives ask whether
a particular form adds anything (\S\ref{sec:method}). New-source tests
are frozen prospectives: forms, coefficients, and predictions are
committed before the compressed model is measured, and a failed
prospective keeps its label.

The evidence has three layers (\S\ref{sec:exp}). On a heterogeneous
panel of 12 public models from five series, the capability response to
compression is strongly capability- and series-specific: QA is the
least-damaged capability in Qwen3 and OLMo-3-32B, where it improves
under moderate pruning, and the most-damaged one in the Gemma series at
mild densities; 3-bit quantization collapses 12 of 14 models but not
OLMo-3-32B. On a controlled Pythia panel, the dense capability vector
is the input that carries prediction, training tokens add little beyond
it, and source-conditioned pruning predictors beat source-free curves
at seen densities and at interpolated densities on the development
sources, mostly for math and code, and match them at a new in-range
size; a shared power form adds nothing over a per-density regression
with a fixed interpolation rule. The same predictors fail at density
$0.55$ and at a source five times larger than the development range,
in both arms. Quantization gains appear
only in the collapse regime and depend on the fitting protocol. For
distillation, a reuse-count coordinate helps QA on an unseen data pool
with a budget-dependent bias, and a constant is best for math and code.
Method-agnostic laws in storage or active-parameter ratios fail a
matched-storage test.

\paragraph{Contributions.}
(1) A capability-conditioned prediction problem with an explicit
admitted-information budget, a common loss interface across three
compression methods, and a two-axis validation design that separates
source transfer from configuration transfer (\S\ref{sec:prelim},
\S\ref{sec:method}).
(2) Candidate relations for each arm, compared on the same inputs
against same-budget baselines, with frozen prospective tests on new
sources and their measured range of validity
(\S\ref{sec:prune}--\ref{sec:distill}).
(3) An empirical account of the conditions under which transferable
prediction is possible and of its failure boundaries, with reproducible
measurements, registers, and evaluation code (\S\ref{sec:synthesis}).
Dense-reference scaling laws are prior work that we use, the Fisher
implementation is infrastructure, and method selection is discussed
but not validated.
